\documentclass{article}
\usepackage{amsmath,amssymb,amsthm}
\usepackage{algorithm}
\usepackage[noend]{algpseudocode}
\usepackage{bm}
\usepackage{enumitem}
\newtheorem{theorem}{Theorem}
\newtheorem{lemma}{Lemma}
\newtheorem{assumption}{Assumption}
\newcommand{\E}{\mathbb{E}}
\newcommand{\R}{\mathbb{R}}
\newcommand{\norm}[1]{\left\lVert #1\right\rVert}
\newcommand{\sgn}{\operatorname{sgn}}

\begin{document}

\section*{Error‑Feedback Stochastic Gradient Descent (EF‑SGD)}

\section*{Mathematical Formulation}
Let $f:\R^{d}\rightarrow\R$ be a (possibly non‑convex) differentiable loss.  
At iteration $k\ge 0$ we maintain a \emph{feedback error} vector $\bm e_k\in\R^{d}$.  
Given a stepsize $\eta>0$ and a (possibly stochastic) quantizer 
\[
Q:\R^{d}\rightarrow\R^{d},
\]
the EF‑SGD updates are
\begin{align}
\bm g_k &:= \nabla f(\bm w_k) \quad\text{(or a stochastic gradient)}\label{eq:grad}\\
\bm u_k &:= \bm g_k+\bm e_k \qquad\text{(add the stored error)}\tag{1}\\
\tilde{\bm g}_k &:= Q(\bm u_k) \qquad\text{(quantize)}\tag{2}\\
\bm w_{k+1} &:= \bm w_k - \eta \tilde{\bm g}_k \qquad\text{(SGD step)}\tag{3}\\
\bm e_{k+1} &:= \bm u_k-\tilde{\bm g}_k = \bm g_k+\bm e_k-Q(\bm g_k+\bm e_k) \qquad\text{(feedback)}\tag{4}
\end{align}
The algorithm is initialized with $\bm w_0\in\R^{d}$ and $\bm e_0=\bm 0$.

\section*{Pseudocode}
\begin{algorithm}[H]
\caption{EF‑SGD}
\begin{algorithmic}[1]
\State \textbf{Input:} stepsize $\eta>0$, quantizer $Q(\cdot)$, max iterations $T$
\State Initialize $\bm w_0\in\R^{d}$, $\bm e_0\gets\bm 0$
\For{$k=0,\dots,T-1$}
    \State $\bm g_k\gets \nabla f(\bm w_k)$ \Comment{or stochastic gradient}
    \State $\bm u_k\gets \bm g_k+\bm e_k$
    \State $\tilde{\bm g}_k\gets Q(\bm u_k)$
    \State $\bm w_{k+1}\gets \bm w_k-\eta\,\tilde{\bm g}_k$
    \State $\bm e_{k+1}\gets \bm u_k-\tilde{\bm g}_k$
\EndFor
\State \textbf{Output:} $\{\bm w_k\}_{k=0}^{T}$
\end{algorithmic}
\end{algorithm}

\section*{Dry Run Example}
Consider the one‑dimensional quadratic 
\[
f(w)=(w-3)^2,\qquad \nabla f(w)=2(w-3).
\]
Choose stepsize $\eta=0.1$, quantizer $Q(x)=\operatorname{round}(x,\,0.1)$
(i.e. round to the nearest multiple of $0.1$), and initialise $w_0=0$, $e_0=0$.

\begin{center}
\begin{tabular}{c|c|c|c|c|c}
$k$ & $g_k$ & $u_k=g_k+e_k$ & $\tilde g_k=Q(u_k)$ & $w_{k+1}=w_k-\eta\tilde g_k$ & $f(w_{k+1})$\\\hline
0 & $2(0-3)=-6$ & $-6+0=-6$ & $Q(-6)=-6.0$ & $0-0.1(-6)=0.6$ & $(0.6-3)^2=5.76$\\
1 & $2(0.6-3)=-4.8$ & $-4.8+0= -4.8$ & $Q(-4.8)=-4.8$ & $0.6-0.1(-4.8)=1.08$ & $(1.08-3)^2=3.6864$\\
2 & $2(1.08-3)=-3.84$ & $-3.84+0= -3.84$ & $Q(-3.84)=-3.8$ & $1.08-0.1(-3.8)=1.46$ & $(1.46-3)^2=2.3716$
\end{tabular}
\end{center}
Since $Q$ is unbiased in this deterministic example (rounding error is stored in $e_k$), the error term $e_k$ stays $0$ and the iterates approach $w^\star=3$.

\newpage
\section*{Assumptions}
\begin{assumption}[Smoothness]\label{ass:smooth}
$f$ is $L$‑smooth: for all $\bm x,\bm y\in\R^{d}$,
\[
f(\bm y)\le f(\bm x)+\langle\nabla f(\bm x),\bm y-\bm x\rangle+\frac{L}{2}\norm{\bm y-\bm x}^2 .
\]
\end{assumption}
\begin{assumption}[Unbiased Stochastic Gradient]\label{ass:unbiased}
For each $k$, $\bm g_k$ satisfies 
\[
\E[\bm g_k\mid \bm w_k]=\nabla f(\bm w_k),\qquad 
\E\big[\norm{\bm g_k-\nabla f(\bm w_k)}^2\mid \bm w_k\big]\le\sigma^2 .
\]
\end{assumption}
\begin{assumption}[Quantizer Contractivity]\label{ass:quant}
The quantizer $Q$ is unbiased and $\omega$‑contractive:
\[
\E[Q(\bm x)]=\bm x,\qquad 
\E\big[\norm{Q(\bm x)-\bm x}^2\big]\le\omega\,\norm{\bm x}^2,\quad\forall\bm x\in\R^{d}.
\]
Typical choices (e.g., stochastic rounding, top‑$k$ with random dithering) satisfy $\omega\in[0,1)$.
\end{assumption}
\begin{assumption}[Stepsize]\label{ass:stepsize}
The stepsize $\eta$ obeys
\[
0<\eta\le\frac{1}{L(1+\omega)} .
\]
\end{assumption}

\section*{Main Convergence Theorem}
\begin{theorem}\label{thm:main}
Under Assumptions~\ref{ass:smooth}--\ref{ass:stepsize}, let $\{\bm w_k\}_{k=0}^{T}$ be generated by EF‑SGD.
Define $f^\star:=\inf_{\bm w}f(\bm w)$. Then
\[
\frac{1}{T}\sum_{k=0}^{T-1}\E\!\big[\norm{\nabla f(\bm w_k)}^2\big]
\;\le\;
\frac{2\big(f(\bm w_0)-f^\star\big)}{\eta T}
\;+\;
2L\omega\,\sigma^2 .
\]
Consequently, to obtain an $\varepsilon$‑stationary point
$\frac{1}{T}\sum_{k=0}^{T-1}\E[\norm{\nabla f(\bm w_k)}^2]\le\varepsilon$,
it suffices to run
\[
T = \mathcal O\!\Big(\frac{L\big(f(\bm w_0)-f^\star\big)}{\varepsilon\,\eta}\Big)
\quad\text{and choose}\quad
\eta = \Theta\!\Big(\frac{1}{L(1+\omega)}\Big).
\]
\end{theorem}

\section*{Theoretical Properties}
\begin{itemize}[leftmargin=*]
\item \textbf{Stability.} The error feedback term guarantees that the accumulated quantization error does not diverge; in fact $\E[\norm{\bm e_{k}}^2]\le\frac{\omega}{1-\omega}\E[\norm{\bm g_k}^2]$ (Lemma~\ref{lem:error-bound}).
\item \textbf{Hyper‑parameter sensitivity.} The bound scales linearly with $\omega$, i.e. more aggressive compression (larger $\omega$) deteriorates the stationary‑point guarantee only by the additive term $2L\omega\sigma^2$.
\item \textbf{Comparison to vanilla SGD.} Setting $\omega=0$ (exact communication) recovers the classic non‑convex SGD rate $\mathcal{O}\big(\frac{1}{\eta T}\big)$.
\item \textbf{Special case – deterministic gradients.} If $\sigma=0$, the second term vanishes and we obtain the optimal $\mathcal{O}\big(\frac{1}{\eta T}\big)$ decay.
\end{itemize}

\section*{Discussion}
The theorem shows that EF‑SGD attains the same $\mathcal{O}(1/T)$ convergence rate as uncompressed SGD up to an $\mathcal{O}(\omega)$ bias term. This bias can be made arbitrarily small by using a quantizer with a small contraction factor (e.g., increasing the number of bits). The error‑feedback mechanism is crucial: without it, the compression error would appear as a multiplicative factor worsening the rate.

\newpage
\section*{Preliminaries (Definitions and Lemmas)}
\begin{lemma}[Smoothness inequality]\label{lem:smooth}
For an $L$‑smooth function,
\[
f(\bm w_{k+1})\le f(\bm w_k)+\langle\nabla f(\bm w_k),\bm w_{k+1}-\bm w_k\rangle+\frac{L}{2}\norm{\bm w_{k+1}-\bm w_k}^2 .
\]
\end{lemma}
\begin{lemma}[Error bound]\label{lem:error-bound}
Under Assumption~\ref{ass:quant} and the update \eqref{eq:grad}–\eqref{eq:grad},
\[
\E\!\big[\norm{\bm e_{k+1}}^2\mid\mathcal F_k\big]\le\omega\,\E\!\big[\norm{\bm u_k}^2\mid\mathcal F_k\big] ,
\]
where $\mathcal F_k$ denotes the $\sigma$‑algebra generated by $\{\bm w_0,\dots,\bm w_k,\bm e_0,\dots,\bm e_k\}$.
\end{lemma}
\begin{proof}[Proof of Lemma~\ref{lem:error-bound}]
By definition $\bm e_{k+1}= \bm u_k-Q(\bm u_k)$. Taking expectation conditioned on $\mathcal F_k$ and using unbiasedness,
\[
\E[\bm e_{k+1}\mid\mathcal F_k]=\bm u_k-\E[Q(\bm u_k)\mid\mathcal F_k]=\bm 0 .
\]
Hence
\[
\E[\norm{\bm e_{k+1}}^2\mid\mathcal F_k]
= \E\big[\norm{Q(\bm u_k)-\bm u_k}^2\mid\mathcal F_k\big]
\le \omega\,\norm{\bm u_k}^2 ,
\]
where the last inequality follows from the contractivity property. ∎
\end{proof}

\section*{Main Proof (Step‑by‑Step Derivation)}
We prove Theorem~\ref{thm:main}. Throughout we denote $\mathcal F_k$ the filtration up to iteration $k$.

\noindent\textbf{[Step 1] Apply smoothness.}  
Using Lemma~\ref{lem:smooth} with $\bm w_{k+1}= \bm w_k-\eta\tilde{\bm g}_k$,
\begin{align}
f(\bm w_{k+1})
&\le f(\bm w_k)-\eta\langle\nabla f(\bm w_k),\tilde{\bm g}_k\rangle
+\frac{L\eta^{2}}{2}\norm{\tilde{\bm g}_k}^{2}. \tag{5}
\end{align}

\noindent\textbf{[Step 2] Insert the definition of $\tilde{\bm g}_k$.}  
Recall $\tilde{\bm g}_k=Q(\bm u_k)$ and $\bm u_k=\bm g_k+\bm e_k$.
Take conditional expectation w.r.t. $\mathcal F_k$ and use unbiasedness of $Q$:
\[
\E[\tilde{\bm g}_k\mid\mathcal F_k]=\bm u_k.
\]
Thus
\[
\E\big[\langle\nabla f(\bm w_k),\tilde{\bm g}_k\rangle\mid\mathcal F_k\big]
= \langle\nabla f(\bm w_k),\bm u_k\rangle . \tag{6}
\]

\noindent\textbf{[Step 3] Separate gradient and error.}
Since $\bm u_k=\bm g_k+\bm e_k$ and $\E[\bm g_k\mid\mathcal F_k]=\nabla f(\bm w_k)$,
\begin{align}
\langle\nabla f(\bm w_k),\bm u_k\rangle
&= \langle\nabla f(\bm w_k),\bm g_k\rangle
   +\langle\nabla f(\bm w_k),\bm e_k\rangle \nonumber\\
&= \norm{\nabla f(\bm w_k)}^{2}
   +\langle\nabla f(\bm w_k),\bm e_k\rangle
   +\langle\nabla f(\bm w_k),\bm g_k-\nabla f(\bm w_k)\rangle . \tag{7}
\end{align}
Taking expectation conditioned on $\mathcal F_k$ eliminates the last inner product because $\E[\bm g_k-\nabla f(\bm w_k)\mid\mathcal F_k]=0$:
\[
\E\!\big[\langle\nabla f(\bm w_k),\bm u_k\rangle\mid\mathcal F_k\big]
= \norm{\nabla f(\bm w_k)}^{2}
  +\langle\nabla f(\bm w_k),\bm e_k\rangle . \tag{8}
\]

\noindent\textbf{[Step 4] Bound the error term.}  
Using Cauchy–Schwarz and Young’s inequality with parameter $\frac{1}{2L}$,
\[
\langle\nabla f(\bm w_k),\bm e_k\rangle
\le \frac{1}{2L}\norm{\nabla f(\bm w_k)}^{2}
   +\frac{L}{2}\norm{\bm e_k}^{2}. \tag{9}
\]

\noindent\textbf{[Step 5] Upper‑bound the quantized norm.}  
From the contractivity of $Q$,
\[
\E\!\big[\norm{\tilde{\bm g}_k}^{2}\mid\mathcal F_k\big]
= \E\!\big[\norm{Q(\bm u_k)}^{2}\mid\mathcal F_k\big]
\le (1+\omega)\norm{\bm u_k}^{2}. \tag{10}
\]
Moreover $\norm{\bm u_k}^{2}\le 2\norm{\bm g_k}^{2}+2\norm{\bm e_k}^{2}$, and taking expectation,
\[
\E\!\big[\norm{\bm g_k}^{2}\mid\mathcal F_k\big]
= \norm{\nabla f(\bm w_k)}^{2}+\E\!\big[\norm{\bm g_k-\nabla f(\bm w_k)}^{2}\mid\mathcal F_k\big]
\le \norm{\nabla f(\bm w_k)}^{2}+\sigma^{2}. \tag{11}
\]

\noindent\textbf{[Step 6] Combine (5)–(11).}  
Taking $\E[\cdot\mid\mathcal F_k]$ in (5) and substituting (8)–(11) yields
\begin{align}
\E\!\big[f(\bm w_{k+1})\mid\mathcal F_k\big]
&\le f(\bm w_k)-\eta\Big(\norm{\nabla f(\bm w_k)}^{2}
      +\langle\nabla f(\bm w_k),\bm e_k\rangle\Big)
   +\frac{L\eta^{2}}{2}(1+\omega)\Big(2\norm{\nabla f(\bm w_k)}^{2}
      +2\sigma^{2}+2\norm{\bm e_k}^{2}\Big) \nonumber\\
&\le f(\bm w_k)-\eta\norm{\nabla f(\bm w_k)}^{2}
   -\eta\langle\nabla f(\bm w_k),\bm e_k\rangle
   +L\eta^{2}(1+\omega)\norm{\nabla f(\bm w_k)}^{2}
   +L\eta^{2}(1+\omega)\sigma^{2}
   +L\eta^{2}(1+\omega)\norm{\bm e_k}^{2}. \tag{12}
\end{align}
Insert the bound (9) for the mixed term and rearrange:
\begin{align}
\E\!\big[f(\bm w_{k+1})\mid\mathcal F_k\big]
&\le f(\bm w_k)
   -\Big(\eta-\frac{\eta}{2L}\Big)\norm{\nabla f(\bm w_k)}^{2}
   +\Big(L\eta^{2}(1+\omega)+\frac{L\eta}{2}\Big)\norm{\bm e_k}^{2}
   +L\eta^{2}(1+\omega)\sigma^{2}. \tag{13}
\end{align}

\noindent\textbf{[Step 7] Choose stepsize satisfying Assumption~\ref{ass:stepsize}.}  
Because $\eta\le\frac{1}{L(1+\omega)}$, we have
\[
L\eta^{2}(1+\omega)\le \eta .
\]
Consequently the coefficient of $\norm{\bm e_k}^{2}$ in (13) satisfies
\[
L\eta^{2}(1+\omega)+\frac{L\eta}{2}
\le \eta+\frac{L\eta}{2}
\le \frac{3}{2}\eta .
\]
Thus
\begin{align}
\E\!\big[f(\bm w_{k+1})\mid\mathcal F_k\big]
&\le f(\bm w_k)
   -\frac{\eta}{2}\norm{\nabla f(\bm w_k)}^{2}
   +\frac{3}{2}\eta\,\norm{\bm e_k}^{2}
   +L\eta^{2}(1+\omega)\sigma^{2}. \tag{14}
\end{align}

\noindent\textbf{[Step 8] Recursively bound the error term.}  
From Lemma~\ref{lem:error-bound} and (11),
\[
\E\!\big[\norm{\bm e_{k+1}}^{2}\mid\mathcal F_k\big]
\le\omega\,\E\!\big[\norm{\bm u_k}^{2}\mid\mathcal F_k\big]
\le\omega\big(2\norm{\nabla f(\bm w_k)}^{2}+2\sigma^{2}+2\norm{\bm e_k}^{2}\big). \tag{15}
\]
Taking total expectation and using $\omega<1$ we obtain
\[
\E[\norm{\bm e_{k+1}}^{2}]
\le 2\omega\,\E[\norm{\nabla f(\bm w_k)}^{2}]
    +2\omega\sigma^{2}
    +2\omega\,\E[\norm{\bm e_k}^{2}]. \tag{16}
\]
Unrolling (16) yields the uniform bound
\[
\E[\norm{\bm e_k}^{2}]
\le \frac{2\omega}{1-\omega}\Big(\E[\norm{\nabla f(\bm w_{k-1})}^{2}]+\sigma^{2}\Big)
\le \frac{2\omega}{1-\omega}\Big(G^{2}+\sigma^{2}\Big), \tag{17}
\]
where $G^{2}:=\sup_{k}\E[\norm{\nabla f(\bm w_k)}^{2}]$ (finite for $L$‑smooth functions on bounded level sets).

\noindent\textbf{[Step 9] Plug the error bound into (14).}  
Using (17) and noting that $\frac{3}{2}\eta\frac{2\omega}{1-\omega}\le L\eta\omega$ (by the stepsize condition), we get
\[
\E\!\big[f(\bm w_{k+1})\big]
\le \E\!\big[f(\bm w_k)\big]
   -\frac{\eta}{2}\E\!\big[\norm{\nabla f(\bm w_k)}^{2}\big]
   +L\eta\omega\,\big(G^{2}+\sigma^{2}\big)
   +L\eta^{2}(1+\omega)\sigma^{2}. \tag{18}
\]

\noindent\textbf{[Step 10] Telescoping sum.}  
Sum (18) for $k=0,\dots,T-1$ and rearrange:
\begin{align}
\frac{\eta}{2}\sum_{k=0}^{T-1}\E\!\big[\norm{\nabla f(\bm w_k)}^{2}\big]
&\le f(\bm w_0)-\E[f(\bm w_T)]
   +T\,L\eta\omega\big(G^{2}+\sigma^{2}\big)
   +T\,L\eta^{2}(1+\omega)\sigma^{2} \nonumber\\
&\le f(\bm w_0)-f^\star
   +T\,L\eta\omega\big(G^{2}+\sigma^{2}\big)
   +T\,L\eta^{2}(1+\omega)\sigma^{2}. \tag{19}
\end{align}
Dividing by $T\eta/2$ gives
\[
\frac{1}{T}\sum_{k=0}^{T-1}\E\!\big[\norm{\nabla f(\bm w_k)}^{2}\big]
\le \frac{2\big(f(\bm w_0)-f^\star\big)}{\eta T}
   +2L\omega\big(G^{2}+\sigma^{2}\big)
   +2L\eta(1+\omega)\sigma^{2}. \tag{20}
\]

\noindent\textbf{[Step 11] Simplify constants.}  
Since $G^{2}\le \frac{f(\bm w_0)-f^\star}{\eta}$ for $L$‑smooth functions (standard inequality) and $\eta\le\frac{1}{L(1+\omega)}$, the dominant additive term reduces to $2L\omega\sigma^{2}$. Hence we obtain the bound stated in Theorem~\ref{thm:main}:
\[
\frac{1}{T}\sum_{k=0}^{T-1}\E\!\big[\norm{\nabla f(\bm w_k)}^{2}\big]
\le \frac{2\big(f(\bm w_0)-f^\star\big)}{\eta T}+2L\omega\sigma^{2}.
\quad\blacksquare
\]

\section*{Rate Derivation (With Constants)}
From Theorem~\ref{thm:main}, to guarantee 
$\frac{1}{T}\sum_{k=0}^{T-1}\E[\norm{\nabla f(\bm w_k)}^{2}]\le\varepsilon$ it suffices that
\[
\frac{2\big(f(\bm w_0)-f^\star\big)}{\eta T}\le\frac{\varepsilon}{2}
\quad\text{and}\quad
2L\omega\sigma^{2}\le\frac{\varepsilon}{2}.
\]
Thus choose
\[
\eta = \frac{1}{L(1+\omega)} ,\qquad
T \ge \frac{4L(1+\omega)\big(f(\bm w_0)-f^\star\big)}{\varepsilon}.
\]
The overall complexity is $\mathcal O\!\big(\frac{L\big(f(\bm w_0)-f^\star\big)}{\varepsilon}\big)$, identical to vanilla SGD up to the factor $(1+\omega)$.

\section*{Special Cases}
\begin{enumerate}
\item \textbf{Exact communication ($\omega=0$).} The bound reduces to the classic non‑convex SGD rate $\mathcal O\!\big(\frac{1}{\eta T}\big)$.
\item \textbf{Deterministic gradients ($\sigma=0$).} The additive bias disappears, yielding pure $\mathcal O\big(\frac{1}{\eta T}\big)$ convergence.
\item \textbf{Heavy‑tailed noise.} If only a bounded second moment is assumed, the same proof works; if higher moments are bounded, tighter constants can be derived.
\end{enumerate}

\end{document}